% BariatricRSD — NeurIPS 2026 skeleton
% Drop this content into the official NeurIPS 2026 LaTeX template (neurips_2026.sty).
% Download template: https://neurips.cc/Conferences/2026
%
% \documentclass{article}
% \usepackage[final]{neurips_2026}    % use [preprint] while writing, remove for camera-ready
% \usepackage{amsmath,amsfonts,booktabs,graphicx,hyperref}

\title{Structured Temporal Prediction under Workflow Variability:\\
Data-Driven Phase-Order Conditioning for Surgical Video}

\author{%
  % NeurIPS anonymous during review; fill for camera-ready.
  Anonymous Author(s)\\
  \texttt{anon@example.com}\\
}

\begin{document}
\maketitle

\begin{abstract}
  Long-horizon prediction on surgical video—remaining surgery duration (RSD), phase recognition,
  and intraoperative deviation detection—is hard because different surgeons execute the same
  procedure in materially different orders. Most existing methods apply generic temporal attention
  and treat every surgery as interchangeable, leaving cross-surgeon and cross-center variability
  unexploited. We propose a simple addition to a standard video Transformer: a learnable
  phase-order conditioning token, computed offline from the procedural signature of each training
  video and prepended to the encoder sequence as a conditioning analog of a \texttt{[CLS]} token.
  On MultiBypass140 (140 Roux-en-Y gastric bypass videos, 2 centers), this conditioning—when
  derived from a data-driven k-means clustering of phase-transition bigrams—reduces within-center
  RSD mean absolute error from \textbf{13.55 to 12.27 min} and halves epoch-level variance from
  0.94 to 0.49~min. A naive heuristic clustering, in contrast, hurts the model. Cross-center
  transfer remains substantially degraded (18.05~min, +5.78~min gap), quantifying the
  workflow-variability problem the method is designed for. The useful contribution is not a new
  architectural block but the general principle that structured temporal prediction under
  heterogeneous workflows benefits from explicit, data-driven workflow conditioning; we provide
  matched-compute ablations supporting this claim.
\end{abstract}

\section{Introduction}
% ~1 page. Problem setup, workflow variability as core challenge, three contributions.
% Content: paste from paper/draft.md §1.

\section{Related Work}
% ~0.75 page. Surgical RSD, MultiBypass140, surgical VLP, conditional tokens, multi-task.
% Content: §2 of draft.md.

\section{Method}
\label{sec:method}
% ~1.5 pages. Include Figure 1 (architecture from notebooks/methods_comparison.ipynb Fig 5).

\subsection{Problem formulation}
% Define inputs, three outputs (RSD, deviation, phase), normalization.

\subsection{Architecture}
\begin{figure}[t]
  \centering
  % \includegraphics[width=\linewidth]{figs/architecture.pdf}
  \caption{BariatricRSD architecture. The phase-order conditioning token (purple) is the novel
    component; all other blocks are adapted from prior work.}
  \label{fig:architecture}
\end{figure}

\subsection{Phase-order clustering}
% Describe: phase bigrams -> TF-IDF -> PCA -> kmeans. Emphasize data-driven vs. heuristic.

\subsection{Multi-task loss}
\begin{equation}
  \mathcal{L}_{\text{total}} = \sum_{k \in \{\text{rsd}, \text{dev}, \text{phase}\}}
    \exp(-\sigma_k)\,\mathcal{L}_k + \sigma_k,
\end{equation}
% with log-variance $\sigma_k$ learned per task (Kendall et al. 2018).

\section{Experiments}
\label{sec:experiments}

\subsection{Dataset and splits}
% MultiBypass140 overview. 5-fold CV, fold 0. Cross-center split: Bern (train) -> Stras (val).

\subsection{Implementation details}
% Hyperparameters. ViT-B/16, HTA (6 blocks), batch 64, AdamW 1e-4, cosine 15ep, GH200.

\subsection{Main result: H1 --- phase-order conditioning}
\label{sec:h1}
\begin{table}[t]
  \centering
  \caption{Within-center H1 ablation (MultiBypass140 fold 0, matched compute, 15 epochs each).
    Phase-order conditioning with data-driven k-means clusters improves every metric.}
  \label{tab:h1}
  \begin{tabular}{lrrr}
    \toprule
    Configuration & Min val MAE (min) & Mean val MAE & Std (epochs) \\
    \midrule
    Without phase-order & 13.55 & 14.33 & 0.94 \\
    With kmeans phase-order (\textbf{ours}) & \textbf{12.27} & \textbf{12.84} & \textbf{0.49} \\
    $\Delta$ (ours $-$ ablation) & $-1.28$ & $-1.49$ & $-0.45$ \\
    \bottomrule
  \end{tabular}
\end{table}

\subsection{H2 --- Multi-task does not hurt RSD}
% Table with Run 001 vs Run 004.

\subsection{Cross-center generalization}
\begin{table}[t]
  \centering
  \caption{Cross-center transfer. A $\sim$6-minute MAE gap quantifies workflow variability
    between Bern and Strasbourg cohorts.}
  \label{tab:crosscenter}
  \begin{tabular}{lrr}
    \toprule
    Metric & Within-center & Cross-center (Bern$\to$Stras) \\
    \midrule
    Min val MAE (min) & 12.27 & \textbf{18.05} \\
    Mean val MAE     & 12.84 & 18.94 \\
    Pearson $r$      & 0.80 & 0.69 \\
    \bottomrule
  \end{tabular}
\end{table}

\subsection{Clustering quality matters}
% Narrative around Run 001 (heuristic, worse) vs Run 006 (kmeans, best). Negative control.

\subsection{Three-seed replication}
% TODO -- pending.

\section{Discussion and Limitations}
% §5 of draft.md. Honest about what we are not claiming.

\section{Conclusion}
% §6 of draft.md.

\section*{Acknowledgments}
% Funding, compute credits, dataset contributors.

\small
\bibliographystyle{plainnat}
\bibliography{references}

\newpage
\appendix

\section{Full hyperparameters per run}
\section{Clustering sensitivity (k = 4, 6, 8)}
\section{Per-epoch validation curves for all runs}
\section{Qualitative case studies}
\section{NeurIPS Paper Checklist}
% Fill per https://neurips.cc/public/guides/PaperChecklist

\end{document}
